% Cloudra V2.0 软件使用说明书（软著申请用），tectonic 直接编译
\documentclass[zihao=-4,a4paper]{ctexart}
\usepackage[top=25mm,bottom=25mm,left=25mm,right=25mm]{geometry}
\usepackage{ctex}
\usepackage{booktabs}
\usepackage{array}
\usepackage{amsmath}
\emergencystretch=2em
\usepackage{enumitem}
\usepackage{fancyhdr}
\usepackage{hyperref}
\hypersetup{hidelinks}
\usepackage{caption}
\captionsetup{labelsep=quad}
\usepackage{xcolor}
\usepackage{listings}
\setmonofont{DejaVu Sans Mono}
\setCJKmonofont{WenQuanYi Zen Hei}
\lstset{
  basicstyle=\ttfamily\small,
  frame=single,
  rulecolor=\color{gray!60},
  breaklines=true,
  keepspaces=true,
  aboveskip=8pt,
  belowskip=8pt,
}
\pagestyle{fancy}
\fancyhf{}
\fancyhead[L]{\small Cloudra V2.0 软件使用说明书}
\fancyhead[R]{\small 第 \thepage{} 页}
\renewcommand{\headrulewidth}{0.4pt}

\title{\bfseries Cloudra 软件使用说明书\\[6pt] \large V2.0}
\author{Cloudra Team}
\date{2026年8月17日}

\begin{document}

\maketitle
\thispagestyle{empty}

\tableofcontents
\newpage

\section{软件概述}

\subsection{软件名称与版本}
\begin{itemize}[leftmargin=4em,itemsep=2pt]
  \item 软件全称：Cloudra化学科研AI Agent协作框架软件
  \item 软件简称：Cloudra
  \item 版本号：V2.0
  \item 开发完成日期：2026年8月17日
\end{itemize}

\subsection{软件用途}
Cloudra 是一款面向化学科研领域的 AI Agent 协作框架软件。它以事件驱动的运行时内核协调一个总控 Agent（Coordinator）与八个专业领域子 Agent，为化学研究人员提供文献检索、分子结构处理、光谱数据解析、量子化学（DFT）与分子动力学（MD）计算任务提交与监控、机理图（DAG）引导的反应机理研究以及科研文档撰写等一站式辅助服务。

\subsection{运行环境}
\begin{itemize}[leftmargin=4em,itemsep=2pt]
  \item 操作系统：Linux、macOS、Windows
  \item 编程语言：Python 3.11 及以上
  \item 依赖环境：OpenAI 兼容的大语言模型 API
  \item 界面环境：Electron 桌面客户端或现代浏览器（Chrome / Edge / Firefox）
\end{itemize}

\section{软件功能模块}

\subsection{事件驱动运行时内核}
软件的核心是事件驱动的运行时内核（runtime\_v2），由以下组件构成：
\begin{itemize}[itemsep=2pt]
  \item \textbf{事件日志与事件总线}：全部会话行为（消息、工具调用、审批、提醒等）以事件形式顺序记录并广播，前端通过 SSE 实时订阅，完整日志可持久化重放。
  \item \textbf{命令总线}：所有写操作先事件后状态，保证多 Agent 并发下的状态一致性；对已失败或取消的 Agent 发送消息时自动重开新一轮运行，用户无需手动重启。
  \item \textbf{Agent 会话状态机}：每个 Agent 拥有独立身份与状态（idle、running、waiting\_approval、waiting\_input、finished 等），重启会话保留身份并基于目标、DAG、长期记忆与工作区外化重建上下文。
  \item \textbf{并发运行器}：每个 Agent 一个异步任务，真正并行执行且可取消；单轮崩溃原地自动重试（默认 2 次、退避 2/4 秒），历史仅记录成功轮次，保证推理过程可恢复。
  \item \textbf{固定轮次管道}：每轮依次经过权限策略、DAG 策略、目标策略三个可插拔阶段；LLM 输出以 token 增量流式推送，每轮完整推理过程落持久化事件。
\end{itemize}

\subsection{多 Agent 协作体系}
\begin{itemize}[itemsep=2pt]
  \item \textbf{Coordinator（协调者）}：负责接收用户需求、拆解任务并调度各专业子 Agent。
  \item \textbf{八类专业子 Agent}：见表~\ref{tab:agents}。
  \item \textbf{子 Agent 调度与蜂群}：支持按模板批量展开的 Swarm 并发调度（可设并发上限）、同步等待与父任务取消级联。
  \item \textbf{AgentHub 论坛}：Agent 间公告板，支持发帖、浏览、检索与跨 Agent 来电呼叫。
  \item \textbf{统一信箱}：用户与 Agent 间消息统一投递，支持定向回复与呼叫。
  \item \textbf{提醒调度}：拦截定时工具调用，以内存定时器在到期时自动投递提醒并唤醒对应 Agent。
\end{itemize}

\begin{table}[htbp]
  \centering
  \caption{专业化 Agent 矩阵}
  \label{tab:agents}
  \begin{tabular}{ll>{\raggedright\arraybackslash}p{8.5cm}}
    \toprule
    Agent 名称 & 标识 & 功能定位 \\
    \midrule
    Coordinator & a-0 & 任务协调与总控 \\
    Explorer & a-1 & 项目文件结构导航与发现 \\
    Writer & a-2 & 文档撰写、报告总结、代码编写 \\
    Researcher & a-3 & 学术文献检索与专业数据调研 \\
    small\_mol & a-4 & 小分子建模、构象与结构文件处理 \\
    md & a-5 & 分子动力学模拟、力场与轨迹分析 \\
    dft & a-6 & 量子化学（Gaussian）输入输出与计算流程 \\
    experiment & a-7 & 实验数据处理与操作辅助 \\
    env & a-8 & 系统环境配置与软件安装管理 \\
    \bottomrule
  \end{tabular}
\end{table}

\subsection{推理与上下文管理}
\begin{itemize}[itemsep=2pt]
  \item \textbf{流式输出}：LLM token 增量以事件流直达前端，逐字渲染；瞬时数据不落日志、不参与重放。
  \item \textbf{持久化推理记录}：每轮完整思考过程落一条持久事件，历史轮次刷新后仍可回看。
  \item \textbf{常驻与一次性注入}：目标、DAG 等常驻说明以固定注入方式参与每轮组装；通知类一次性注入在轮内统一收集。
  \item \textbf{手动上下文压缩}：对超长对话可一键压缩历史（非流式摘要 + 工具配对边界对齐），在轮边界安全执行。
  \item \textbf{长期记忆}：MEMORY.md 长期记忆文件自动加载，跨会话保留研究偏好与结论。
\end{itemize}

\subsection{工具套件}
\begin{itemize}[itemsep=2pt]
  \item \textbf{Chemtool 化学工具}：SMILES 生成与转换（smiles\_to\_mol）、Gaussian 输入文件编写（gjf\_writer）、量子化学输出解析（SCF 能量、振动频率、过渡态虚频判别等）、NMR/IR 光谱处理与格式转换。
  \item \textbf{Searchtool 检索工具}：arXiv、OpenAlex（含 ChemRxiv 聚合）学术文献检索，原生联网搜索适配。
  \item \textbf{Bashtool 命令工具}：跨平台 Shell 执行，内置命令分级、文件变更审批与危险命令守卫。
  \item \textbf{Writetool 写作工具}：结构化文档、实验记录与报告写入。
  \item \textbf{Plantool 计划工具}：研究计划与机理图（DAG）操作，提供机理草稿、证据绑定等十余种图操作指令。
  \item \textbf{Managetool 管理工具}：子 Agent 调度、蜂群展开、AgentHub 论坛、定时提醒。
  \item \textbf{Skilltool 技能工具}：加载科研方法技能包（DFT 计算、文献综述、论文写作等）。
  \item \textbf{Stoptool 终止工具}：按目标完成度门禁终止会话，防止过早收尾。
\end{itemize}

\subsection{DAG 机理图引导编排}
软件支持按会话开启的 DAG 引导编排模式，用于有机反应机理研究：
\begin{itemize}[itemsep=2pt]
  \item \textbf{机理草稿与阶段契约}：模型提出机理草稿（节点与基元步骤边），经接受后落为交付图；阶段按契约顺序绑定。
  \item \textbf{调用复核}：声明为证据源的工具调用完成后自动进入复核收件箱，由机制检查器（如过渡态虚频数 $n_{\text{imag}}=1$）确定性派生证据等级。
  \item \textbf{覆盖率审计}：停止会话时审计交付覆盖率，存在未决缺口时阻止提交，保证机理闭环。
  \item \textbf{督导模式}：可选的科学督导 Agent 接受/驳回草稿、逐调用复核并绑定阶段。
\end{itemize}

\subsection{安全与权限系统}
\begin{itemize}[itemsep=2pt]
  \item \textbf{五级权限模型}：READ（读）、WRITE（写）、MANAGE（管理子 Agent）、SCRIPT（执行脚本）、DESTRUCTIVE（破坏性操作）。
  \item \textbf{危险操作审批流}：危险工具调用挂起并等待用户确认；请求与决策全部事件化、可追溯。
  \item \textbf{YOLO 自动审批守卫}：可选全自动模式，由独立 LLM 守卫对审批请求进行安全判断。
  \item \textbf{专项守卫}：坐标守卫（拦截绕过审查直接改写量子化学坐标文件的写法）、命令习惯守卫（提示 heredoc/sed -i 等易错写文件命令）、委派守卫（提示 Coordinator 自行提交计算任务）。
\end{itemize}

\subsection{计算执行与数据管理}
\begin{itemize}[itemsep=2pt]
  \item \textbf{本地与 Slurm 计算}：计算代理扫描会话库认领就绪任务，本地拉起监视进程持续解析输出；支持多节点注册表与 Slurm 作业提交、轮询与续租。
  \item \textbf{SQLite 持久化}：会话级数据库承载事件日志、DAG 图存储与复核收件箱，支持 WAL 模式与增量重放恢复。
  \item \textbf{不可变产物快照}：工具产生的文件以内容寻址快照保存，支持短期下载授权。
  \item \textbf{备份恢复}：内置 backup 子命令，在线快照 + SHA-256 清单，支持校验与恢复。
\end{itemize}

\section{软件安装与配置}

\subsection{安装步骤}
\begin{lstlisting}[language=bash]
# 克隆代码仓库
git clone <仓库地址>
cd <仓库目录>

# 安装（依赖统一声明在 pyproject.toml；可选安装化学扩展）
pip install -e .
pip install -e ".[chem]"   # 可选：RDKit、nmrglue 等化学依赖
\end{lstlisting}

\subsection{初始化配置}
\begin{lstlisting}[language=bash]
cloudra config init    # 交互式初始化：LLM 地址、模型、密钥等
cloudra config show    # 查看当前配置（敏感信息自动脱敏）
cloudra config set coordinator.model gpt-4o   # 设置单个配置项
cloudra config get coordinator.model           # 读取单个配置项
\end{lstlisting}

\section{软件使用方法}

\subsection{启动图形界面}
\begin{lstlisting}[language=bash]
cloudra                 # 构建并启动 Electron 桌面前端
cloudra web --browser   # 在浏览器中打开前端（单端口模式）
\end{lstlisting}

前端界面由三部分组成：
\begin{itemize}[itemsep=2pt]
  \item \textbf{项目栏与 Agent 栏}：切换会话项目、查看各 Agent 状态（空闲/运行/等待审批/等待输入）。
  \item \textbf{对话区}：流式对话输出、推理过程分段折叠、工具调用行、Markdown 与 KaTeX 公式渲染。
  \item \textbf{检查器（Inspector）}：六个标签页，管理目标、DAG 编排模式、权限、审批请求、提醒与上下文压缩。
\end{itemize}

\subsection{常用管理指令}
\begin{lstlisting}[language=bash]
cloudra --version      # 查看版本
cloudra logs <对话ID>  # 查看某个对话的日志
cloudra skills         # 列出可用的科研技能包
cloudra backup create  # 创建数据备份（含 SHA-256 清单）
cloudra backup verify  # 校验备份完整性
cloudra backup restore # 恢复到空目录
\end{lstlisting}

\section{典型使用场景}

\subsection{反应机理研究（DAG 模式）}
用户在项目设置中开启 DAG 引导编排，向 Coordinator 提出目标（如“研究某催化循环的完整机理”）。系统引导模型提出机理草稿，经接受后按阶段契约绑定证据：DFT 子 Agent 提交过渡态计算，计算作业完成后自动进入复核收件箱，由虚频检查器确定证据等级；全部必需阶段达成后经覆盖率审计交付完整机理图。

\subsection{DFT 计算全流程}
用户描述目标分子与计算需求，Coordinator 调度 dft 子 Agent：先生成 Gaussian 输入文件（经坐标守卫校验来源），提交计算后由监视进程实时解析输出，自动提取 SCF 能量、HOMO--LUMO 能隙、振动频率等关键指标并写入报告。

\subsection{文献调研与综述}
用户要求调研某课题最新进展，Researcher 子 Agent 检索 arXiv 与 OpenAlex 数据库，Writer 子 Agent 依据检索结果撰写综述，产出结构与引用规范的 Markdown 报告。

\section{技术支持}
\begin{itemize}[leftmargin=4em,itemsep=2pt]
  \item 开发团队：Cloudra Team
  \item 反馈渠道：项目仓库 Issue 页面
\end{itemize}

\vfill
\begin{center}
\small 本说明书适用于 Cloudra V2.0 版本软件
\end{center}

\end{document}
